\documentclass[10.5pt,a4paper,twoside,twocolumn]{ctexart}

%============== 页面设置 ==============
\usepackage[top=2.54cm,bottom=2.54cm,left=2.2cm,right=2.2cm,
  headheight=13pt,headsep=3mm]{geometry}
\usepackage{graphicx}
\usepackage{subfig} 
\usepackage{url}
\usepackage{amsmath,amsthm}
\usepackage{amssymb,amsfonts}
\usepackage{booktabs}
\usepackage{float}
\usepackage{caption}
\usepackage{fancyhdr}
\usepackage{xcolor}
\usepackage[numbers]{natbib}
\usepackage[hidelinks]{hyperref}


%============== 字体 ==============
\setmainfont{Times New Roman}

%============== 页眉 ==============
\pagestyle{fancy}
\fancyhf{}
\fancyhead[CO]{\small 基于分层博弈的ICVPS最优安全控制技术报告}
\fancyhead[RO]{\small 顾作一 2515100292}
\fancyhead[LE]{\small 顾作一 2515100292}
\fancyhead[CE]{\small 基于分层博弈的ICVPS最优安全控制技术报告}
\renewcommand{\headrulewidth}{0.4pt}
\setlength{\headwidth}{\textwidth}

%============== 自定义命令 ==============
\newcommand{\T}{\mathsf{T}}

%============== 标题 ==============
\title{\heiti\zihao{2} 基于分层博弈的智能网联车队列系统\\最优安全控制技术报告}
\author{\fangsong\zihao{3} 顾作一 2515100292\\\zihao{6}（宁波大学，浙江 宁波 315000）}
\date{}

\begin{document}

\twocolumn[
\maketitle

\vspace{-6mm}
\noindent\rule{\textwidth}{0.4pt}

\vspace{2mm}
{\heiti\zihao{5} 摘\quad 要}\quad
\zihao{5}
针对智能网联车队列系统（ICVPS）在恶意虚假数据注入（FDI）攻击下面临的安全控制挑战，本文从博弈论与最优控制理论的双重视角，对Zhiwei Xia等人提出的一种基于分层博弈的最优安全控制方法进行深入的理论分析。该方法构建了融合外层微分图博弈与内层零和动态对抗博弈的分层框架：内层通过求解Hamilton-Jacobi-Isaacs（HJI）方程的鞍点获得车辆级最优安全控制策略与最坏攻击策略的解析形式；外层将每辆车建模为图智能体，驱动整个ICVPS达到纳什均衡，确保全局队列稳定性。为克服HJI方程的解析求解困难，论文设计了改进余弦退火机制的离线学习残差神经网络（OLRNN）进行数值逼近。本报告从系统建模、博弈理论、最优控制和机器学习四个维度，对该方法的数学模型、博弈框架的理论基础、纳什均衡的证明逻辑以及OLRNN的训练理论进行了系统解析。

\vspace{3mm}
\noindent{\heiti\zihao{5} 关键词}\quad \zihao{5} 智能网联车队列系统；虚假数据注入攻击；分层博弈安全控制；微分图博弈；零和博弈；HJI方程；纳什均衡；离线学习残差神经网络

\vspace{3mm}
\begin{center}
\zihao{3}\textbf{Hierarchical Game-Based Optimal Security Control of Vehicular Platoon Systems: A Technical Report}\\
\vspace{2mm}
\zihao{5} GU Zuo-Yi 2515100292\\
\zihao{6} (Ningbo University, Ningbo 315000, China)
\end{center}

\zihao{5}\noindent\textbf{Abstract}\quad This report provides a comprehensive theoretical analysis of a hierarchical game-based optimal security control framework for Intelligent Connected Vehicle Platoon Systems (ICVPS) under malicious False Data Injection (FDI) attacks. The framework integrates outer-layer differential graphical games with inner-layer zero-sum dynamic adversarial games. At the vehicle control layer, the saddle point of the HJI equation yields optimal security control policies; at the network communication layer, the differential graphical game drives the multi-vehicle system to Nash equilibrium. An offline-learning residual neural network with improved cosine annealing is employed to solve the HJI equation numerically. The analysis covers system modeling, game theory, optimal control, and machine learning.

\vspace{3mm}
\noindent\textbf{Keywords}\quad vehicular platoon systems; false data injection attacks; hierarchical game security control; differential graphical game; zero-sum game; HJI equation; Nash equilibrium; offline learning residual neural network

\vspace{5mm}
\noindent\rule{\textwidth}{0.4pt}
\vspace{5mm}
]

%====================================================================
\section{\heiti\zihao{4} 引言}
%====================================================================

随着人工智能、5G通信与自动驾驶技术的深度融合，智能网联车（Intelligent Connected Vehicles, ICVs）已成为现代交通运输系统转型的核心引擎。智能网联车队列系统（Intelligent Connected Vehicle Platoon Systems, ICVPS）作为车联网（V2X）的典型应用场景，通过使一组车辆以预设间距和速度保持协调编队行驶，在提升道路通行效率、降低能耗、改善驾驶体验方面展现出巨大潜力。该技术被广泛视为未来智能交通系统（ITS）的关键使能技术\cite{jia2016}。

然而，无线通信网络的开放性使ICVPS面临严重的网络攻击威胁。虚假数据注入（False Data Injection, FDI）攻击通过篡改传感器测量数据或控制指令，使控制器基于虚假信息运行，危及队列稳定性与行车安全。与传统FDI攻击不同，恶意FDI攻击具备智能性和自适应性——其攻击强度与时机可依据系统的实时状态动态调整，从而具有更强的隐蔽性和破坏性\cite{petit2015,geng2025}。这一特性使得传统的基于固定阈值的检测与防御方法难以奏效。目前，针对ICVPS的FDI攻击防御方法主要包括基于状态估计的检测方法\cite{sargolzaei2021}、事件触发机制\cite{yang2022}以及自适应控制方案\cite{ren2020}等，但是这些方法往往忽视了攻击与系统性能之间的动态博弈关系，导致性能下降或防御成本过高。

为应对这一挑战，上海大学的Xia等人在IEEE Transactions on Vehicular Technology（2026年第75卷第6期）发表了一项开创性研究\cite{xia2026}，提出了一个分层博弈最优安全控制（OSC）框架。该框架的核心理论创新在于：将博弈论与最优控制理论有机融合，构建了"博弈的博弈"体系——在内层，通过零和动态对抗博弈模型刻画控制器与恶意攻击之间的对抗关系，将安全控制问题转化为HJI方程的鞍点求解问题；在外层，运用微分图博弈理论驱动多车系统达到纳什均衡，实现全局协同的最优化。求解层面，论文提出了改进余弦退火机制的离线学习残差神经网络（OLRNN），以数值方式逼近最优解。

本技术报告从系统建模、博弈理论、最优控制和机器学习四个维度，对该方法的理论基础进行深入分析，并结合交通运输行业的特点，评估其在真实场景中的工程可行性与应用前景。

%====================================================================
\section{\heiti\zihao{4} 系统建模与理论分析}
%====================================================================

\subsection{\heiti\zihao{5} 车辆动力学模型与误差系统}

论文考虑由$n+1$辆智能车辆组成的ICVPS，其中车辆0为虚拟领航车，车辆$1$至$n$为跟随车。第$i$辆跟随车的动力学由三阶微分方程描述：
\begin{equation}
\begin{cases}
	\dot{p}_i(t) = v_i(t) \\
	\dot{v}_i(t) = a_i(t) \\
	\dot{a}_i(t) = K_i(t) + \frac{1}{\tau_i} u_i(t) + f_{a(i,-i)}(t) \\
\hfill +  w_i(t) + \sum_{j \in N_i} \pi_{ij}    	f_{a(j,-j)}(t) \\
	y_i(t) = p_i(t) + f_{y(i,-i)}(t) + \sum_{j \in N_i} \pi_{ij} f_{y(j,-j)}(t)
\end{cases}
\end{equation}
位置$p_i(t)$、速度$v_i(t)$、加速度$a_i(t)$为三阶积分链，$u_i, y_i, \tau_i$ 为控制输入、系统输出和发动机时间滞后；并引入了发动机时滞$\tau_i$、未建模动态分量$K_i(t)$、外部扰动$w_i(t)$以及分别作用于控制输入通道和传感器输出通道的FDI攻击$f_a$和$f_y$。

论文引入了网络影响系数$\pi_{ij}$来描述攻击在通信网络中的传播效应：攻击不仅直接作用于目标车辆$j$，还通过传播项$\sum_{j\in\mathcal{N}_i} \pi_{ij} f_{a(j,-j)}(t)$影响到其邻居车辆$i$。这一建模精确刻画了V2V通信拓扑中攻击的级联扩散特性——在交通运输场景中具有重要的物理意义：高速公路上队列中一辆车被攻破后，虚假信息会通过网络拓扑级联传播，引发连锁反应。

安全约束要求每辆车与前车间保持$i\cdot d$的安全距离，同时匹配领航车的速度和加速度。定义跟踪误差向量：

\begin{equation}
e_i = \left[ y_i - (y_0 - i \cdot d),\; v_i - v_0,\; a_i - a_0 \right]^{\T}
\label{eq:error}
\end{equation}

基于该误差定义，论文推导出误差动态方程：

\begin{equation}
\dot{e}_i = q_i(e_i) + m_i u_i + f_i(e_i, e_{-i}) + \sum_{j\in\mathcal{N}_i} \pi_{ij} f_j(e_j, e_{-j})
\label{eq:wcdt}
\end{equation}

\subsection{\heiti\zihao{5} 控制目标}
论文的控制目标是在恶意FDI攻击和外部扰动同时存在的条件下，设计最优安全控制策略 $u_i$ 和相应的最坏情况攻击策略 $f_i$，使得：
（1）\textbf{确保车辆队列系统的跟踪稳定性：}这意味着队列中的车辆需保持均匀的车距，并具有允许的安全裕度 $\alpha d$；同时，其速度和加速度必须与虚拟领航车保持一致。该目标的数学表达如下：

\begin{equation}
	\begin{aligned}
	& \lim_{t \to \infty} \|y_i - y_0 + id\| \le \alpha d \\
	& \lim_{t \to \infty} \|v_i(t) - v_0(t)\| = 0 \\
	& \lim_{t \to \infty} \|a_i(t) - a_0(t)\| = 0 \quad
    \end{aligned}
\end{equation}
其中，$\alpha$ 为安全参数（用于界定防碰撞的安全距离阈值）。

（2）\textbf{实现局部车辆性能与全局队列性能的同步优化：}在恶意攻击和干扰下，保证每辆车局部最优的同时维持全局最优队列编队。

\subsection{\heiti\zihao{5} 恶意FDI攻击的智能特性建模}

论文对FDI攻击的建模是其理论贡献之一。恶意FDI攻击具备两个智能特性。

（1）\textbf{状态依赖的自适应调整}。攻击信号根据车辆自身的跟踪误差$e_i$及其邻居的跟踪误差$e_{-i}$动态调整：

\begin{equation}
f_{a(i,-i)} = S_i \left[ e_i^{\T}, e_{-i}^{\T} \right]^{\T},\quad
f_{y(i,-i)} = \Upsilon_i \left[ e_i^{\T}, e_{-i}^{\T} \right]^{\T}
\label{eq:attack}
\end{equation}

其中$S_i(t)$和$\Upsilon_i(t)$为时变攻击增益矩阵。当系统偏离理想状态越远（即误差越大），攻击的幅度也越大，形成正反馈式的恶性循环。

（2）\textbf{攻击时机的选择性触发}。攻击在车辆状态发生显著波动时激活——攻击者能够识别系统最脆弱的时刻，例如在领航车加减速的过渡阶段，此时控制器的调节负担最重、系统抗扰动能力最弱，从而最大化其破坏效果。

论文采用最坏情况假设（worst-case assumption），即攻击者能够完全获取车辆的实时状态信息。这一假设给出了系统脆弱性的上界，据此设计的防御策略在最恶劣的攻击条件下仍能保证系统性能，从而在实际（通常弱于该假设）的攻击环境中具有更强的鲁棒性。

\subsection{\heiti\zihao{5} 代价函数与控制目标的形式化}

论文将安全控制问题形式化为最优控制问题。定义性能指标：

\begin{equation}
\begin{aligned}
J_i = \int_0^{\infty} \Big[ & e_i^{\T} Q_i e_i + u_i^{\T} R_{ii} u_i \\
& + \sum_{j \in \mathcal{N}_i} u_j^{\T} R_{ij} u_j - \sum_{j \in \mathcal{N}_i} \theta_j f_j^{\T} C_{ij} f_j \Big] dt
\end{aligned}
\label{eq:cost}
\end{equation}

该代价函数具有明确的理论物理意义：（1）$e_i^{\T} Q_i e_i$项衡量系统的跟踪精度；（2）$u_i^{\T} R_{ii} u_i$项惩罚控制能量消耗；（3）$\sum u_j^{\T} R_{ij} u_j$项反映了邻居车辆控制输入对当前车辆性能的影响，这是图博弈区别于独立优化问题的关键所在；（4）负向攻击项$-\sum \theta_j f_j^{\T} C_{ij} f_j$体现了零和博弈的本质：攻击者试图最大化$J_i$，而控制器试图最小化$J_i$，在鞍点处达到均衡。

%====================================================================
\section{\heiti\zihao{4} 核心成果：分层博弈安全控制理论框架}
%====================================================================

该论文核心的理论创新在于提出了"博弈的博弈"（game of games）分层框架（图\ref{fig:SOC}）。这一框架通过两层博弈的有机结合，将全局队列协同与局部车辆安全两个层面的优化问题统一处理。

\begin{figure}[htbp]
	\centering
	\includegraphics[width=0.5\textwidth]{图2.png}
	\caption{分层博弈最优安全控制（OSC）框架}
	\label{fig:SOC}
\end{figure}

\subsection{\heiti\zihao{5} 内层零和动态对抗博弈}

在车辆控制层，控制器与恶意FDI攻击构成典型的二人零和动态博弈。Xu等人\cite{xu2022}在互联系统的分布式最优容错控制中采用了Stackelberg微分图博弈方法，与之不同的是，本文分析的方法采用同时行动的零和博弈框架，根据 \textbf{Nash-Pontryagin极小极大原理}，只要存在鞍点，该零和博弈就存在唯一解，满足以下等式：

\begin{equation}
	\begin{aligned}
\min_{u_i} \max_{f_i} J_i(\bar{e}_i, u_i, & u_j^*, f_i, f_j^*) =  \\
& \max_{f_i} \min_{u_i} J_i(\bar{e}_i, u_i, u_j^*, f_i, f_j^*) \quad
\end{aligned}
\end{equation}

其中 $u_i^*$ 为车辆 $i$ 的最优控制策略，$f_i^*$ 为影响车辆 $i$ 的最坏情况攻击策略。



其求解可转化为Hamilton-Jacobi-Isaacs（HJI）方程。定义车辆 $i$ 的哈密顿（Hamiltonian）函数为：


\begin{equation}
\begin{aligned}
	H_i&(\bar{e}_i, \nabla J_i, u_i, u_j^*, f_i, f_j) = \\
	&\bar{e}_i^T Q_i \bar{e}_i + u_i^T E_{ii} u_i - \theta_i f_i^T C_{ii} f_i - \sum_{j \in N_i} \theta_j f_j^T C_{ij} f_j \\
	&+ \nabla J_i^T \left( q_i + m_i u_i + f_i + \sum_{j \in N_i} \pi_{ij} f_j \right)
\end{aligned}
\end{equation}

其中 $\nabla J_i$ 表示 $J_i$ 的梯度。根据最优性条件 $0 = \min_{u_i} \max_{f_i} H_i$，可求得\textbf{最优控制策略 $u_i^*$ }和\textbf{最坏情况攻击策略 $f_i^*$}：

\begin{equation}
	\begin{aligned}
	u_i^* = -\frac{1}{2} E_{ii}^{-1} m_i^T \nabla J_i^*, \quad f_i^* = \frac{1}{2\theta_i} C_{ii}^{-1} \nabla J_i^* \quad	
	\end{aligned}
\end{equation}

最优性能指标$\mathcal{J}_i^*$分解为车辆自身代价$\mathcal{J}_{ii}^*$与邻车耦合代价$\sum_{j\in N_i}\mathcal{J}_{ij}^*$，最优控制$u_i^*$、最坏攻击策略$f_i^*$可进一步展开并代入哈密顿函数，并令其等于 0，可推导出 HJI 偏微分方程：

\begin{equation}
\begin{aligned}
	\bar{e}_i^T Q_i \bar{e}_i &- \frac{1}{4} (\nabla J_i^*)^T m_i E_{ii}^{-1} m_i^T \nabla J_i^* \\
	&+ \frac{1}{4\theta_i} (\nabla J_i^*)^T C_{ii}^{-1} \nabla J_i^* + (\nabla J_i^*)^T q_i(e_i) \\
	&- \sum_{j \in N_i} \frac{1}{2\theta_j} (\nabla J_i^*)^T \pi_{ij} C_{jj}^{-1} \nabla J_j^* \\
	&- \sum_{j \in N_i} \frac{1}{4\theta_j} (\nabla J_j^*)^T C_{jj}^{-1} C_{ij} C_{jj}^{-1} \nabla J_j^* = 0 \quad 
\end{aligned}
\end{equation}

将$\mathcal{J}_i^*(0)=0$代入，得到假设3.1：存在正常数 $a_i, b_i, c_i, d_i$，使得最优性能指标满足：

\begin{equation}
	\begin{aligned}
 a_i \|\bar{e}_i\|^2 \le J_i^*(\bar{e}_i) \le b_i \|\bar{e}_i\|^2, \\ 
 c_i \|\bar{e}_i\| \le \|\nabla J_i^*(\bar{e}_i)\| \le d_i \|\bar{e}_i\|
	\end{aligned}
	\label{eq:xingneng}
\end{equation}

推得定理3.1：在假设 2.1、2.2、3.1 成立前提下，采用式 \eqref{eq:xingneng}最优安全控制器，若满足以下两个条件，则车队系统误差动力学式\eqref{eq:wcdt} 渐近稳定：

1.  对于每辆车 $i$，参数 $\theta_i$ 满足攻击能量阈值条件：

\begin{equation}
	\begin{aligned}
		\theta_i \ge \theta_{i,\min} \triangleq \frac{d_i^2 \|C_{ii}^{-1}\|}{4\lambda_{\min}(Q_i) + \frac{c_i^2}{4 m_i \lambda_{\min}(E_{ii}^{-1}) m_i}} 
	\end{aligned}
\end{equation}

2.  对于ICVPS中的任意闭环路径 $i_1 \to i_2 \to \cdots \to i_w \to i_1$，互联增益满足\textbf{小增益定理}条件：

\begin{equation}
	\begin{aligned}
	\gamma_{i_1 i_2} \gamma_{i_2 i_3} \cdots \gamma_{i_w i_1} < 1
	\end{aligned}
\end{equation}

其中，从 $V_j$ 到 $V_i$ 的增益定义为：$\gamma_{ij} \triangleq \frac{h_i \bar{n}_i z_{ij}}{a_j k_i}$，且 $z_{ij} \triangleq \frac{d_j^2}{4\theta_j} \|C_{jj}^{-1} C_{ij} C_{jj}^{-1}\|$。

注释3.1：定理 3.1基于2个条件，条件 1 要求攻击 Lipschitz 上界\(l_{i}(\zeta_{i}) \leq l_{i, max }(\zeta_{i})\)，最优安全控制领域普遍采用攻击有界假设，贴合工程实际。条件 2 要求车队拓扑所有环路增益乘积小于1；本文车队拓扑为单环结构，所有车辆控制器参数同步整定即可满足约束。因此条件成立时，定理 3.1说明分层博弈最优安全控制器具备智能抗攻击能力，可自适应恶意 FDI 攻击，生成最小代价最优控制策略。

\subsection{\heiti\zihao{5} 外层微分图博弈理论}

在外层网络通信层，论文将多车队列系统建模为微分图博弈。设图$G = (\mathcal{N}, \mathcal{E})$，其中$\mathcal{N} = \{1, 2, \dots, n\}$为车辆节点集合，$\mathcal{E}$为通信边集合。每辆车的性能指标$J_i$不仅依赖于自身策略$(u_i, f_i)$，还依赖于邻居的策略$(u_j, f_j)$，这种耦合性质正是图博弈的核心特征。与Wang等人\cite{wang2024}针对USV系统的Stackelberg博弈方法不同，本框架采用完全非合作的微分图博弈模型，使所有车辆在平等地位下达到纳什均衡。

纳什均衡的精确定义为：若存在$n$-元策略对集合$\{(u_1^*, f_1^*), (u_2^*, f_2^*), \dots, (u_n^*, f_n^*)\}$使得：

\begin{equation}
\begin{cases}
	J_i(u_i^*, u_j^*, f_i^*, f_j^*) \le J_i(u_i, u_j^*, f_i^*, f_j^*) \\
	J_i(u_i^*, u_j^*, f_i^*, f_j^*) \ge J_i(u_i^*, u_j^*, f_i, f_j^*)
\end{cases}
\label{eq:nash}
\end{equation}

则该策略集合构成纳什均衡。左侧不等式表示：若车辆$i$的攻击策略$f_i$偏离最优$f_i^*$，其性能指标反而降低，因此攻击者没有动机偏离；右侧表示：若控制策略$u_i$偏离最优$u_i^*$，性能指标上升，因此控制器也没有动机偏离。

\textbf{定理3.2：}对于ICVPS，上述策略集构成\textbf{纳什均衡}。在该均衡下，每辆车的性能指标达到最优值，即：

\begin{equation}
	\begin{aligned}
		J_i^*(\bar{e}_i) = V_i(\bar{e}_i), \quad \forall i \in \mathcal{N}
	\end{aligned}
	\label{eq:zb}
\end{equation}

注释 3.2：定理 3.2 说明单车最优性能同时依赖自身与邻车策略；纳什均衡具备严格最优性，任何单车单方面改变策略均无法提升自身性能，且会降低车队整体协同效率。本文分层博弈框架不存在局部、全局性能权衡损耗，可在恶意 FDI 攻击与外部扰动下同时实现车队全局最优协同与单车局部稳定最优。

注释 3.3：分层博弈架构大幅降低大规模车队计算复杂度：全局问题分层拆解，上层负责编队规划、生成基准策略，下层单车仅求解邻域局部对抗子问题。各子问题状态维度固定，HJI 方程求解复杂度不随车辆总数$n$指数增长；单车控制决策可并行独立计算，仅需自身与邻车状态信息。

论文中两层博弈之间的耦合机制：内层零和博弈的输出（车辆$i$的最优控制策略$u_i^*$）构成了外层图博弈中车辆$i$的行动集；而外层图博弈的代价耦合（通过$R_{ij}$和$C_{ij}$实现）又反过来影响内层HJI方程中值函数的求解。这种"博弈嵌套博弈"的设计避免了将全部问题纳入单一大型博弈所导致的"维数灾难"问题。

由于分层设计将全局问题分解为局部子问题，每个子问题仅涉及有限的状态维度（车辆自身及其直接邻居的状态），HJI方程求解的复杂度不会随车辆总数$n$呈指数增长，而是保持为常数。每辆车的控制决策可以并行独立计算，仅需自身及邻居的信息——这一性质称为可扩展性（scalability），是算法从理论走向工程实践的关键前提。Marelli和Fu\cite{marelli2023}提出的网络增益定理为分析这类互联系统的稳定性提供了理论工具，保证了分层设计的数学严谨性。

%====================================================================
\subsection{\heiti\zihao{5} 离线学习残差神经网络的理论分析}
%====================================================================

由于HJI方程的解析解仅在极简条件下存在（如线性二次型问题），对于含有非线性动力学和耦合博弈的ICVPS而言，必须寻求数值逼近方法。论文选用残差神经网络（ResNet）作为函数逼近器来近似值函数$V_i(\bar{e}_i)$，进而通过式\ref{eq:zb}得到近似最优控制策略。

\subsubsection{\heiti\zihao{5} 残差网络的理论优势}

残差网络的核心设计是引入跳跃连接（skip connection），每个残差块实现变换：

\begin{equation}
h_{\text{out}} = h_{\text{in}} + \sigma(W h_{\text{in}} + b)
\label{eq:resnet}
\end{equation}
\(h_{in}/h_{out}\)为残差块输入输出；\(W、b\)为权重与偏置；\(\sigma\)为激活函数。

从函数逼近理论角度看，残差结构具有两个重要优势：第一，缓解了深层网络中的梯度消失问题——由于恒等映射的存在，梯度可以直接通过跳跃连接反向传播到浅层，使得深层网络的训练更加稳定；第二，残差结构实际上是在学习状态的增量（即残差）而非完整的映射——网络只需要在恒等映射基础上进行"修正"，而非从头学习复杂的函数关系，从而显著降低了学习难度和所需样本量。

\subsubsection{\heiti\zihao{5} 改进余弦退火机制的理论分析}

论文提出了一种改进的余弦退火算法，其核心创新在于动态确定最大学习率，而非使用固定的预定义值。该算法包含两个阶段。

\textbf{阶段1：学习率范围测试（LR Range Test）}。算法以指数增长的方式扫描学习率：

\begin{equation}
\eta_{tk} = \eta_{\text{start}} \times \left( \frac{\eta_{\text{end}}}{\eta_{\text{start}}} \right)^{\frac{tk}{Tk}}
\label{eq:lrscan}
\end{equation}

在每次测试迭代中记录损失$L_{tk}$，建立$\log(\eta)$-Loss关系曲线，损失开始急剧上升的点对应的学习率即为$\eta_{\max}(tk)$。这一方法基于学习率与损失函数Lipschitz常数的关系：当学习率超过损失函数梯度Lipschitz常数的倒数时，梯度下降不再收敛。

\textbf{阶段2：改进的余弦退火调度}。将动态$\eta_{\max}(tk)$代入余弦退火公式：

\begin{equation}
\eta_{tk} = \eta_{\min} + \frac{1}{2} (\eta_{\max}(tk) - \eta_{\min}) \left( 1 + \cos\left( \frac{Tk^{\text{cur}}}{Tk^{\max}} \pi \right) \right)
\label{eq:cosine}
\end{equation}

\(\eta_{min}\)为最小学习率；\(Tk_{cur}\)为当前周期迭代步数；\(Tk_{max}\)为单周期总迭代步数。

改进版本使得$\eta_{\max}$能够根据损失曲面的局部曲率自适应调整。当损失曲面平坦时（低曲率），$\eta_{\max}$可以较大以加快收敛；当损失曲面陡峭时（高曲率），$\eta_{\max}$自动减小以防止发散。这种自适应机制使得网络在训练过程中能够持续逃离局部最优，同时保持稳定收敛。

\subsubsection{\heiti\zihao{5} HJI残差最小化的损失函数}

OLRNN的训练通过最小化HJI残差来实现。损失函数定义为：

\begin{equation}
	\begin{aligned}
	\mathcal{L}(\theta) & = \mathbb{E}_{\bar{e}_i \sim \rho(\bar{e}_i)} \\
	& \left[ \mathcal{L}_0(\bar{e}_i, u^*, f^*) + \nabla V_{NN}(\bar{e}_i; \theta)^{\T} D(\bar{e}_i, u^*, f^*) \right]^2
	\end{aligned}
	\label{eq:loss}
\end{equation}

其中$\rho(\bar{e}_i)$是状态样本的分布，$D(\cdot)$包含系统的动力学项，$u^*$和$f^*$依赖于$\nabla V_{NN}$。从理论上看，当损失趋近于零时，$V_{NN}$趋近于真实的值函数$V_i$，从而保证了近似最优策略逼近真实最优策略。论文分析了计算复杂度：训练深度神经网络的时间复杂度与训练样本数、网络参数规模和迭代次数成正比，保持在多项式级别，显著低于HJI方程的直接求解。所有训练在云端高性能计算平台上离线完成，在线阶段仅需加载训练好的网络参数，实现快速实时计算。

%====================================================================
\section{\heiti\zihao{4} 仿真验证与结果分析}
%====================================================================

论文基于CARLA高保真仿真平台进行验证。实验配置包括1辆虚拟领航车和4辆跟随车。非线性车辆动力学模型为：

\begin{equation}
K_i = -(a_i + l_i v_i^2 M_i^{-1} + o_i M_i^{-1})\tau_i^{-1} - 2l_i v_i a_i M_i^{-1} 
\end{equation}

其中，$l_i, o_i$ 分别表示空气阻力系数和机械阻力。外部干扰设定为 $w_i(t)= 0.1 \sin(t)$，控制输入受限，满足 $\|u_i(t)\| \le 150 \text{ N}\cdot\text{m}$。纵向参数选择为 $l_i= 0.34$，$o_i= 0.1 \text{ N}$，期望车间距 $d= 10 \text{ m}$，安全系数 $\alpha= 0.2$。其他相关参数列于\ref{tab:car}中 [28]。

\begin{table}[htbp]
	\centering
	\caption{车队车辆参数}
	\label{tab:car}
	\begin{tabular}{|c|c|c|c|c|}
		\hline
		车辆编号$i$ & 1 & 2 & 3 & 4 \\
		\hline
		空气阻力系数\(l_i\)(m) & 3.5 & 4.5 & 5 & 5 \\
		发动机时滞\(\tau_i\)(s) & 0.4 & 0.85 & 0.5 & 0.55 \\
		整车质量\(M_i(\times10^3\ \text{kg})\) & 1.8 & 1.85 & 1.8 & 1.9 \\
		\hline
	\end{tabular}
\end{table}

虚拟领航车（$i=0$）的加速度定义如下：

\begin{equation}
a_0(t)= 
\begin{cases} 
	0.48t \text{ m/s}^2 & 1 \text{ s} \le t \le 5 \text{ s} \\ 
	2.4 \text{ m/s}^2 & 5 \text{ s} \le t \le 10 \text{ s} \\ 
	-0.48t+ 7.2 \text{ m/s}^2 & 10 \text{ s} \le t \le 15 \text{ s} \\ 
	0 \text{ m/s}^2 & \text{其他} 
\end{cases}
\end{equation}

仿真中设计了两类攻击以验证所提方法：（1）恶意FDI攻击（9-11s减速阶段、14-16s过渡阶段），攻击增益$S_i(t) = 0.2\cos(t) I_{1 \times 6}$，$\Upsilon_i(t) = 0.1\sin(2t) I_{1 \times 6}$，在速度变化的关键时刻触发，持续2秒每周期；（2）传统FDI攻击（29-31s匀速阶段），作为对比基准。前两次恶意攻击造成的干扰显著强于第三次传统攻击，验证了恶意FDI攻击相较于传统FDI攻击的更大威胁性。

与基于融合状态估计补偿器的自适应控制律\cite{pan2024}对比，所提方法在以下方面具有显著优势：（1）收敛速度更快（图\ref{fig:compare}）——攻击结束后的恢复时间更短；（2）最大跟踪误差更小——加速度、速度和位置误差的峰值均显著降低；（3）控制目标可达性（图\ref{fig:gz_compare}）——所提方法能使所有车辆的位置偏差保持在安全距离阈值$\alpha d$以内，而基准方法无法实现这一控制目标。CARLA仿真截图（图\ref{fig:fz}）直观地展示了分层博弈OSC算法使车辆平稳快速地收敛到参考轨迹。

\begin{figure*}[htbp]
	\centering
	\subfloat[论文方法]{\includegraphics[width=0.45\textwidth]{图3.png}\label{fig:traj_proposed}}
	\hfill
	\subfloat[基准方法]{\includegraphics[width=0.45\textwidth]{图4.png}\label{fig:traj_baseline}}
	\caption{不同方法车辆轨迹对比}
	\label{fig:compare}
\end{figure*}

\begin{figure*}[htbp]
	\centering
	\subfloat[论文方法]{\includegraphics[width=0.45\textwidth]{图5.png}\label{fig:gz_proposed}}
	\hfill
	\subfloat[基准方法]{\includegraphics[width=0.45\textwidth]{图6.png}\label{fig:gz_baseline}}
	\caption{不同方法车辆轨迹对比}
	\label{fig:gz_compare}
\end{figure*}

\begin{figure}[htbp]
	\centering
	\includegraphics[width=0.5\textwidth]{图7.png}
	\caption{车队仿真图}
	\label{fig:fz}
\end{figure}

\begin{figure}[htbp]
	\centering
	\includegraphics[width=0.5\textwidth]{图8.png}
	\caption{不同学习率测试}
	\label{fig:xxl}
\end{figure}

\begin{figure}[htbp]
	\centering
	\includegraphics[width=0.5\textwidth]{图9.png}
	\caption{损失函数降低对比}
	\label{fig:ss_compare}
\end{figure}


消融实验对比了三种训练方案：论文提出的改进余弦退火方法、恒定学习率基线和传统余弦退火（图\ref{fig:fz}、\ref{fig:ss_compare}）。结果表明改进方法在加速损失下降方面显著优于其他方案，且在后续优化阶段保持持续有效性。这些消融实验从计算学习理论的角度证明了自适应学习率调度对于解决HJI方程这类高度非线性函数的神经网络逼近问题至关重要。

%====================================================================
\section{\heiti\zihao{4} 理论贡献总结与工程应用评估}
%====================================================================

\subsection{\heiti\zihao{5} 主要理论贡献}

该论文的理论贡献可归纳为四个层面：

（1）\textbf{博弈理论层面}：首次将微分图博弈与零和动态对抗博弈有机融合，构建了面向ICVPS的安全控制"博弈的博弈"框架。在博弈论视角下统一了"个体最优"（内层零和博弈的鞍点解）与"群体最优"（外层图博弈的纳什均衡），为多智能体系统的安全控制问题提供了新的理论范式。

（2）\textbf{最优控制理论层面}：通过HJI方程的理论推导，将带有恶意FDI攻击的非线性ICVPS控制问题转化为动态博弈的鞍点求解问题。给出了最优控制策略和最坏攻击策略的解析形式，并利用Lyapunov理论证明了系统的渐近稳定性。

（3）\textbf{计算学习理论层面}：设计了改进余弦退火机制的离线学习残差神经网络，通过动态学习率调度有效克服了传统方法在训练中陷入局部最优的问题。离线学习避免了在线学习在实际部署中的安全风险。

（4）\textbf{可扩展性分析层面}：论证了分层框架的计算复杂度不随车辆数量指数增长，保持了多项式级别的复杂度，且每辆车的控制决策可以并行独立计算。

\subsection{\heiti\zihao{5} 工程应用评估}

从交通运输工程视角评估，该研究的工程可行性体现在：（1）\textbf{车-云协同架构适配}——OLRNN在云端完成离线训练，车辆端仅需加载网络参数进行前向推理，与当前主流的车-云协同架构高度兼容；（2）\textbf{分布式部署可行}——每辆车的控制决策仅需自身及邻居的状态信息，不依赖全局信息，能适应V2X通信的范围限制和拓扑变化；（3）\textbf{安全冗余设计}——最坏情况假设保证了算法在最恶劣条件下的鲁棒性，在实际中提供充分的安全裕量。

论文也承认了当前方法的局限性：仅考虑单一类型的FDI攻击而未覆盖混合攻击（如DoS+FDI联合攻击）；离线学习的控制性能在持续运行中无法自我改进。这些方向构成了后续研究的重要课题。

%====================================================================
\section{\heiti\zihao{4} 结论}
%====================================================================

本技术报告从系统建模、博弈理论、最优控制和机器学习四个理论维度，对Xia等人提出的基于分层博弈的ICVPS最优安全控制方法进行了深入的理论分析。该方法通过内层零和对抗博弈实现了车辆级的最优安全控制，通过外层微分图博弈实现了队列级的全局协调，并利用改进余弦退火机制的离线学习残差神经网络数值求解了HJI方程。论文在纳什均衡存在性、系统渐近稳定性、计算复杂度可控性等方面提供了严格的理论证明，形成了较为完善的理论体系。

从交通运输行业视角来看，该研究为智能网联车队列系统的网络安全防护提供了坚实的理论基础和可行的工程路径。随着车联网和自动驾驶技术的持续演进，该分层博弈安全控制框架有望在未来的智能交通系统中发挥重要作用。后续研究应聚焦于混合攻击场景下的理论拓展、离线-在线协同学习机制的建立，以及真实道路环境中的部署验证。

\vspace{5mm}


\bibliographystyle{gbt7714-numeric}
\bibliography{ref}

\end{document}
